\documentclass[11pt,a4paper]{article}
\usepackage[utf8]{inputenc}
\usepackage[vietnamese]{babel}
\usepackage{amsmath,amssymb}
\usepackage[version=4]{mhchem}
\usepackage{geometry}
\geometry{top=1.5cm,bottom=1.2cm,left=1.4cm,right=1.4cm,headheight=28pt,headsep=12pt,footskip=18pt}
\usepackage{graphicx}
\usepackage{xcolor}
\usepackage{tabularx}
\usepackage{array}
\usepackage{tikz}
\usepackage{fancyhdr}
\usepackage{enumitem}
\usepackage{multido}
\usepackage{tcolorbox}
\usepackage{colortbl}

\definecolor{hfRed}{HTML}{D32F2F}
\definecolor{hfGreen}{HTML}{8BC390}
\definecolor{mainGreen}{HTML}{2E7D32}
\definecolor{solBlue}{HTML}{1565C0}
\definecolor{tableGreen}{HTML}{E8F5E9}

\usepackage{ex_test}

\pagestyle{fancy}
\fancyhf{}
\renewcommand{\headrulewidth}{0.75pt}
\renewcommand{\headrule}{\hbox to\headwidth{\color{hfGreen}\leaders\hrule height \headrulewidth\hfill}}
\renewcommand{\footrulewidth}{0.75pt}
\renewcommand{\footrule}{\hbox to\headwidth{\color{hfGreen}\leaders\hrule height \footrulewidth\hfill}}

\fancyhead[L]{\small\itshape\color{hfRed}Tài liệu lưu hành nội bộ}
\fancyhead[R]{\small\color{hfRed}Thầy \textbf{TRẦN VĂN THẠNH} - 0777.470.803}

\fancyfoot[L]{\footnotesize\itshape\color{hfRed}CS1: 6/15 Nguyễn Hoàng, P. Kim Long, TP Huế \quad|\quad CS2: 24 Đặng Thái Thân, TP Huế}
\fancyfoot[R]{\small\textbf{\color{hfRed}Trang \thepage}}

\setlength{\parindent}{0pt}
\setlength{\parskip}{2pt}

\begin{document}

\noindent
\begin{tabularx}{\textwidth}{@{}X r@{}}
\textbf{TRƯỜNG THPT CHUYÊN / THPT TỈNH} & \textbf{ĐỀ KIỂM TRA GIỮA HỌC KÌ I} \\
\textbf{ĐỀ CHÍNH THỨC} & \textbf{MÔN: HOÁ HỌC - LỚP 11} \\
\textit{(Đề thi có 04 trang)} & \textbf{Thời gian làm bài: 50 phút} \textit{(không kể thời gian phát đề)} \\
& \textbf{MÃ ĐỀ THI: 102}
\end{tabularx}

\vspace{2pt}
\noindent\textbf{Họ, tên thí sinh:} \dotfill \textbf{Lớp:} \dots\dots\dots \textbf{Số báo danh:} \dots\dots\dots

\vspace{2pt}
\noindent\textit{(Cho biết nguyên tử khối: Cu = 64; Fe = 56; N = 14; O = 16; S = 32; H = 1; K = 39; Cl = 35,5; C = 12; Ca = 40).}
\vspace{2pt}
\hrule height 0.5pt
\vspace{6pt}

\noindent{\color{mainGreen}\textbf{\Large A. PHẦN TRẮC NGHIỆM (7,0 ĐIỂM)}}
\vspace{4pt}

\noindent{\color{mainGreen}\textbf{\large PHẦN I. Câu trắc nghiệm nhiều phương án lựa chọn.}} \\
\textit{Thí sinh trả lời từ câu 1 đến câu 12. Mỗi câu hỏi thí sinh chỉ chọn một phương án. (3,0 điểm)}
\vspace{4pt}
\begin{ex}
Sự chuyển dịch cân bằng hóa học là sự dịch chuyển từ trạng thái cân bằng hóa học này
\choice
{sang trạng thái không cân bằng do tác động của các yếu tố từ bên ngoài tác động lên cân bằng.}
{sang trạng thái cân bằng hóa học khác do tác động của các yếu tố từ bên ngoài tác động lên cân bằng.}
{sang trạng thái cân bằng hóa học khác không cần có tác động của các yếu tố từ bên ngoài tác động lên cân bằng.}
{sang trạng thái cân bằng hóa học khác do cân bằng hóa học tác động lên các yếu tố bên ngoài.}
\loigiai{}
\end{ex}
\begin{ex}
Cho \ce{Cu} phản ứng với dung dịch \ce{HNO3} loãng, nóng thu được một chất khí không màu hóa nâu trong không khí, khí đó là
\choice
{\ce{NO}.}
{\ce{NH3}.}
{\ce{N2O}.}
{\ce{N2}.}
\loigiai{}
\end{ex}
\begin{ex}
Trong thực tế, khi bảo quản dung dịch muối sắt(III) chloride (\ce{FeCl3}) trong phòng thí nghiệm, người ta thường phải cho thêm vài giọt dung dịch \ce{HCl} đặc vào lọ đựng. Việc làm này nhằm mục đích chủ yếu là:
\choice
{Ngăn cản quá trình bay hơi của dung dịch \ce{FeCl3}, làm cho dung dịch trong suốt, không màu.}
{Ức chế sự thủy phân của ion \ce{Fe^3+} tạo kết tủa \ce{Fe(OH)3} và ngăn sự oxi hóa \ce{Fe^2+} (nếu có).}
{Làm tăng màu vàng nâu đặc trưng của dung dịch \ce{FeCl3}.}
{Tăng tốc độ điện li của muối \ce{FeCl3}.}
\loigiai{}
\end{ex}
\begin{ex}
Ở điều kiện thường, sulfur tồn tại ở dạng tinh thể, được tạo nên từ các phân tử sulfur. Số nguyên tử trong mỗi phân tử sulfur là
\choice
{4.}
{8.}
{6.}
{2.}
\loigiai{}
\end{ex}
\begin{ex}
Với xúc tác của các ion kim loại trong khói bụi, các oxide của sulfur và nitrogen bị oxi hóa bởi oxygen, ozone, hydrogen peroxide, gốc tự do,... rồi hòa tan vào nước tạo thành các acid tương ứng. Hai acid tạo thành từ quá trình trên là
\choice
{\ce{H2S} và \ce{H2SO4}.}
{\ce{H2CO3} và \ce{HNO3}.}
{\ce{H2S} và \ce{HNO3}.}
{\ce{H2SO4} và \ce{HNO3}.}
\loigiai{}
\end{ex}
\begin{ex}
Trong khí quyển, nguyên tố nitrogen tồn tại chủ yếu dưới dạng
\choice
{\ce{N2}.}
{\ce{NH4+}.}
{\ce{NO3-}.}
{\ce{NO2-}.}
\loigiai{}
\end{ex}
\begin{ex}
Hiện tượng phú dưỡng có thể quan sát được thông qua sự xuất hiện dày đặc của tảo xanh trong nước. Khi đó, lượng oxygen trong nước sẽ nhanh chóng giảm đi. Điều này là do:\\(a) Quá trình phân hủy tảo chết bởi vi khuẩn đã tiêu tốn lượng lớn oxygen trong nước.\\(b) Sự phát triển của tảo xanh đã ngăn cản ánh sáng và không khí chứa oxygen khuếch tán vào nước.\\(c) Sự hoạt động của lượng lớn vi khuẩn đã hấp thụ đáng kể oxygen hòa tan trong nước.\\Số nhận định \textbf{đúng} là
\choice
{1.}
{0.}
{2.}
{3.}
\loigiai{}
\end{ex}
\begin{ex}
Chất nào sau đây thuộc loại chất điện li yếu?
\choice
{\ce{NH4Cl}.}
{\ce{KCl}.}
{\ce{HNO3}.}
{\ce{HF}.}
\loigiai{}
\end{ex}
\begin{ex}
Khí nitrogen ít tan trong nước là do
\choice
{nitrogen có độ âm điện lớn.}
{phân tử nitrogen không phân cực.}
{nitrogen có bán kính nguyên tử nhỏ.}
{phân tử nitrogen có liên kết ba bền vững.}
\loigiai{}
\end{ex}
\begin{ex}
Phân tử ammonia có dạng hình học nào sau đây?
\choice
{Chóp tam giác.}
{Chóp tứ giác.}
{Chữ T.}
{Tam giác đều.}
\loigiai{}
\end{ex}
\begin{ex}
Chất chỉ thị acid - base là những chất
\choice
{Chỉ có khả năng đổi màu trong môi trường base mạnh.}
{Có màu sắc thay đổi tùy thuộc vào giá trị pH của dung dịch.}
{Chỉ có khả năng đổi màu trong môi trường acid mạnh.}
{Có màu sắc không thay đổi khi tiếp xúc với các acid hoặc base.}
\loigiai{}
\end{ex}
\begin{ex}
Khi tiến hành các thí nghiệm khảo sát tính háo nước và tính oxi hóa mạnh của \ce{H2SO4} đặc trên các mẫu vật liệu: tờ giấy trắng, hạt gạo, và đường mía; hiện tượng chung quan sát được trên các mẫu vật này là gì và giải thích \textbf{đúng nhất} là?
\choice
{Cả ba mẫu vật chỉ bị mất nước tạm thời, khi cho rửa lại bằng nước sạch thì các mẫu vật sẽ trở lại trạng thái ban đầu.}
{Giấy và đường mía bị hóa đen do bị cháy do nhiệt độ cao, còn gạo không có hiện tượng gì.}
{Tất cả các mẫu vật đều bị hóa đen; nguyên nhân là do \ce{H2SO4} đặc chiếm giữ các nguyên tố H và O (thành phần cấu tạo nên carbohydrate) để lại khung carbon.}
{H2SO4 đặc thủy phân hoàn toàn các chất này thành dung dịch acid hữu cơ làm tờ giấy trắng, hạt gạo, và đường mía sủi bọt khí hydrogen H2.}
\loigiai{}
\end{ex}

\vspace{6pt}
\noindent{\color{mainGreen}\textbf{\large PHẦN II. Câu trắc nghiệm đúng sai.}} \\
\textit{Thí sinh trả lời từ câu 13 đến câu 14. Trong mỗi ý a), b), c), d) ở mỗi câu, thí sinh chọn đúng hoặc sai. (2,0 điểm)}
\vspace{4pt}

\begin{ex}
Cho phản ứng: $\ce{N2(g) + 3H2(g) <=> 2NH3(g)} \quad \Delta_r H^0_{298} = -92\text{ kJ}$.
\choiceTF
\True Ở một nhiệt độ xác định, biểu thức xác định hằng số cân bằng $K_c = \frac{[\ce{NH3}]^2}{[\ce{N2}][\ce{H2}]^3}$.
Vì phản ứng thuận là phản ứng tỏa nhiệt, nên để nâng cao hiệu suất tổng hợp \ce{NH3} trong công nghiệp, người ta thường thực hiện phản ứng ở nhiệt độ vừa phải $500^\circ\text{C}$.
\ce{NH3} một chất khí không màu, không mùi, làm xanh giấy quỳ tím ẩm.
\True Biến thiên enthalpy của phản ứng có giá trị âm $\Delta_r H^0_{298} < 0$ chứng tỏ hệ tỏa nhiệt ra môi trường khi phản ứng thuận xảy ra.
\loigiai{}
\end{ex}

\begin{ex}
Sulfur dioxide (\ce{SO2}) là một trong các khí gây ô nhiễm môi trường, là nguyên nhân chính gây ra mưa acid và nhiều tác hại đối với sức khỏe con người. Trong công nghiệp, khí này được tạo ra từ phản ứng đốt cháy sulfur theo cân bằng sau: $\ce{S(s) + O2(g) ->[t^o] SO2(g)}$.
\choiceTF
Nước mưa bình thường không ô nhiễm có pH khoảng 5,6 do có \ce{CO2} hòa tan. Nếu nước mưa tại vùng công nghiệp nhiễm \ce{SO2}, tạo ra acid \ce{H2SO4} làm pH giảm xuống dưới 5,0 thì quỳ tím sẽ chuyển sang màu xanh.
\True Trong tự nhiên, nguyên tố sulfur (lưu huỳnh) tồn tại phổ biến ở dạng đơn chất trong lòng đất (mỏ lưu huỳnh) và dạng hợp chất trong các khoáng vật như quặng pyrite và thạch cao.
Để giảm thiểu tác hại của \ce{SO2}, các nhà máy được khuyến khích xây dựng các ống khói thật cao để khí \ce{SO2} khuếch tán nhanh lên tầng cao khí quyển, chuyển hóa triệt để thành chất an toàn mà không gây mưa acid.
\True Để xác định hàm lượng \ce{SO2} trong không khí, người ta dẫn lượng khí này qua dung dịch \ce{KMnO4}. Một học sinh tiến hành chuẩn độ 50 mL dung dịch acid \ce{H2SO4} (được tạo ra sau khi thu khí \ce{SO2} bằng dung dịch \ce{NaOH} 0,0001 M). Thể tích \ce{NaOH} tiêu tốn trung bình là 15 mL khi dung dịch chuyển sang màu hồng nhạt (chỉ thị phenolphthalein). Nồng độ mol của dung dịch \ce{H2SO4} ban đầu thu được là $1{,}5\cdot 10^{-5}\text{ M}$.
\loigiai{}
\end{ex}

\vspace{6pt}
\noindent{\color{mainGreen}\textbf{\large PHẦN III. Câu hỏi trắc nghiệm yêu cầu trả lời ngắn.}} \\
\textit{Thí sinh trả lời từ câu 15 đến câu 18. (2,0 điểm)}
\vspace{4pt}

\begin{ex}
Giá trị pH của dung dịch \ce{H2SO4} 0,05 M là bao nhiêu?
\shortans{1}
\loigiai{
\textbf{Trả lời: 1}\\
Dung dịch \ce{H2SO4} phân li hoàn toàn: $\ce{H2SO4 -> 2H+ + SO4^2-}$.\\
$[\ce{H+}] = 2 \times 0{,}05 = 0{,}1\text{ M} = 10^{-1}\text{ M} \Rightarrow \text{pH} = -\log[\ce{H+}] = 1$.
}
\end{ex}

\begin{ex}
Cho a gam \ce{Cu} phản ứng hoàn toàn với dung dịch \ce{H2SO4} đặc, nóng, dư thu được 24,79 lít khí \ce{SO2} (ở đkc, sản phẩm khử duy nhất). Giá trị của a? \textit{(Làm tròn kết quả đến hàng phần mười)}
\shortans{64,0}
\loigiai{
\textbf{Trả lời: 64,0}\\
Ở đkc: $n_{\ce{SO2}} = \frac{24{,}79}{24{,}79} = 1{,}0\text{ mol}$.\\
Phương trình: $\ce{Cu + 2H2SO4\text{ (đặc)} ->[t^o] CuSO4 + SO2\uparrow + 2H2O}$.\\
Theo phương trình: $n_{\ce{Cu}} = n_{\ce{SO2}} = 1{,}0\text{ mol} \Rightarrow a = 1{,}0 \times 64 = 64{,}0\text{ gam}$.
}
\end{ex}

\begin{ex}
Cho các chất: \ce{Au}, \ce{Al}, \ce{FeO}, \ce{Fe(OH)2}, \ce{CaCO3}, \ce{K2SO4} lần lượt tác dụng với dung dịch \ce{HNO3} loãng. Có bao nhiêu trường hợp xảy ra phản ứng oxi hóa -- khử?
\shortans{3}
\loigiai{
\textbf{Trả lời: 3}\\
Các trường hợp xảy ra phản ứng oxi hóa -- khử gồm:
\begin{itemize}[noitemsep,topsep=1pt]
\item \ce{Al} là kim loại có tính khử: $\ce{Al + 4HNO3 -> Al(NO3)3 + NO\uparrow + 2H2O}$.
\item \ce{FeO} chứa $\ce{Fe^2+}$: $3\ce{FeO} + 10\ce{HNO3} \to 3\ce{Fe(NO3)3} + \ce{NO}\uparrow + 5\ce{H2O}$.
\item \ce{Fe(OH)2} chứa $\ce{Fe^2+}$: $3\ce{Fe(OH)2} + 10\ce{HNO3} \to 3\ce{Fe(NO3)3} + \ce{NO}\uparrow + 8\ce{H2O}$.
\end{itemize}
(\ce{Au} không phản ứng; \ce{CaCO3} là phản ứng trao đổi; \ce{K2SO4} không phản ứng).
}
\end{ex}

\begin{ex}
Xét trạng thái tự nhiên của sulfur:
\begin{enumerate}[label=\arabic*),noitemsep,topsep=2pt]
\item Trong tự nhiên sulfur có cả ở dạng đơn chất và hợp chất.
\item Trong cơ thể người, sulfur có trong thành phần của nhiều protein và enzyme.
\item Dạng hợp chất sulfur thường có ở các khoáng vật như pyrite (\ce{FeS2}), chu sa (\ce{HgS}), ...
\item Dạng đơn chất sulfur thường có ở gần các hồ nước mặn.
\end{enumerate}
Hãy sắp xếp các mệnh đề \textbf{đúng} theo số thứ tự tăng dần (ví dụ 12, 123, ...)
\shortans{123}
\loigiai{
\textbf{Trả lời: 123}\\
Mệnh đề (1), (2), (3) đúng. Mệnh đề (4) sai vì đơn chất sulfur thường có ở vùng mỏ ngầm hoặc gần miệng núi lửa, hồ nước mặn chứa các muối sulfate hòa tan. Thứ tự các mệnh đề đúng: 123.
}
\end{ex}

\vspace{6pt}
\noindent{\color{mainGreen}\textbf{\Large B. PHẦN TỰ LUẬN (3,0 ĐIỂM)}}
\vspace{4pt}

\noindent\textbf{Câu 19. (1,0 điểm)} Trình bày cách pha loãng dung dịch sulfuric acid đặc an toàn và nguyên tắc xử lí sơ bộ khi bị bỏng sulfuric acid.

\vspace{10pt}
\multido{\i=1+1}{6}{\noindent\dotfill\\[4pt]}

\vspace{6pt}
\noindent\textbf{Câu 20. (1,0 điểm)} Cho 100 mL dung dịch \ce{KOH} x M vào dung dịch \ce{(NH4)2SO4} dư, đun nóng nhẹ thu được 2,479 lít khí (đo ở điều kiện chuẩn đkc).\\
a) Khí thu được có thể nhận biết bằng cách nào?\\
b) Tính x?

\vspace{10pt}
\multido{\i=1+1}{6}{\noindent\dotfill\\[4pt]}

\vspace{6pt}
\noindent\textbf{Câu 21. (1,0 điểm)} Cho cân bằng hóa học sau trong pha khí:\\
\centerline{$2\ce{SO2(g) + O2(g) <=> 2SO3(g)} \quad \Delta_r H^0_{298} < 0$}\\
Cân bằng trên dịch chuyển theo chiều nào khi thay đổi các điều kiện sau (giải thích ngắn gọn)?\\
a) Tăng nhiệt độ.\\
b) Tăng áp suất chung của hệ phản ứng.\\
c) Dùng thêm chất xúc tác \ce{V2O5}.\\
d) Tăng nồng độ \ce{SO2}.

\vspace{10pt}
\multido{\i=1+1}{7}{\noindent\dotfill\\[4pt]}

\begin{center}
\textbf{--- HẾT ---}
\end{center}
\end{document}
